\runninghead{Rend the Heavens and Come Down}
\properlane{rend-the-heavens}{Rend the Heavens and Come Down: The Cry Answered and Still Prayed}
\label{sec:rend-the-heavens}

The first question is whether the prophet's cry has been answered, and if
so, how the Church can still pray it. Jerome finds the answer in the past
tense that follows the cry. Thomas Aquinas, commenting on the same verses
eight centuries later, keeps Jerome's answer as one sense of the verb and
prefers another as its literal sense. Augustine, preaching on the psalm from
which the Alleluia verse is taken, says why the Church should still pray
``Give us thy Christ'' when Christ has been given, and Cyril of Jerusalem
teaches that Christ comes twice. The cry is a prayer that God has heard in
the Incarnation and that the Church still prays toward Christ's return.

\subsection{Jerome: the promise kept}

Jerome's commentary on Isaiah treats the whole of chapters 63 and 64 as the prayer of a
people made up of the just and of sinners. When he reaches the cry, he
explains its wish by what the speakers had been given until then, which was
only promises: \latin{Utinam disrumperes caelos, et descenderes; et qui semper
promitteris, tu aliquando promissa compleres}, ``O that thou wouldst rend the
heavens and come down, and that thou, who art always being promised, wouldst
at last fulfil what was promised.'' He places the prayer in time: it was
said before the Saviour came and took our nature from the Virgin's womb
(\work{Commentarii in Isaiam} XVII, PL~24, col.~645). The mountains that melt
before the Lord are for him the adverse powers.

Two verses later Jerome read in the Hebrew what the Latin of the reading
still has: \latin{Descendisti}, ``thou didst come down''. The speakers who
prayed ``O that thou wouldst come down'' now say that he came, and Jerome
takes the change from wish to statement as the sign of a prayer granted:
\latin{Pulchreque juxta Hebraicum, quasi exauditi essent}, ``beautifully,
according to the Hebrew, as if they had been heard''. What came down, he
says, is the Word made flesh, Emmanuel, God with us (col.~646).

One verse later the Douay text says, ``the eye hath not seen, O God, besides
thee, what things thou hast prepared for them that wait for thee'' (64:3;
Douay 64:4). Jerome observes that Paul takes this verse for his own
in the letter read today: ``as it is written: That eye hath not seen, nor ear
heard \ldots\ what things God hath prepared for them that love him'' (1~Cor
2:9), and he insists that Paul gives the sense of Isaiah and does not quote an
apocryphal book (cols.~646--647).

\subsection{Aquinas: three comings in one verb}

Aquinas's commentary on Isaiah is a literal exposition, written early in his
teaching, and on these verses it takes up Jerome's reading and adds to it.
The word \latin{dirumperes}, ``that thou wouldst rend'', is a figure spoken to
one who dwells above the heavens; or it means \latin{exinanires temetipsum,
despiciens majestatem, et carnem assumeres}, ``that thou wouldst empty
thyself, setting thy majesty aside, and take flesh''. He supports it with a
psalm that asks the same thing: ``Lord, bow down thy heavens and descend''
(Ps 144 (143):5). The language of emptying is Paul's to the Philippians (2:7),
though Aquinas does not cite the place.

On \latin{Descendisti} he gives three readings. It may be said of the Lord's
coming down to judge us, \latin{praeteritum pro futuro}, the past tense put
for the future. It may be said
of his coming down into the Virgin's womb. Or, \latin{melius, secundum
litteralem sensum}, ``better, according to the literal sense'', it is a proof
drawn from the past: if you came down now, these things would happen as they
happened when you came down to free your people from Egypt, when, as another
psalm says, ``He bowed the heavens, and came down'' (Ps 18 (17):10)
(\work{Expositio super Isaiam} c.~64). The prophet's people remember the
Exodus and ask for it again.

Aquinas adds one more reading in the same chapter. Of the next verse,
\latin{Occurristi laetanti}, ``Thou hast met him that rejoiceth'' (64:4;
Douay 64:5), he says that God goes to meet in help the one who is still
asking or still on the way to him, and then, \latin{mystice}, names Simeon,
the just and devout man of Jerusalem who was ``waiting for the consolation of
Israel'' (Lk 2:25).

\subsection{Augustine: the vine and the Son of man}

The psalm asks God to come in another figure, as a shepherd. Augustine, preaching on
it, said at the outset what it is about: ``The song here is of the Advent of
the Lord and of our Saviour Jesus Christ, and of His vineyard''; it confesses
``both Christ and the vineyard; that is, Head and Body, King and people,
Shepherd and flock'' (\work{Enarrationes in Psalmos} 79.1).
On the first words, addressed to the Shepherd of Israel, he says, ``He is
being invoked to come, He is being expected until He come, He is being
yearned for until He come'' (79.2). On the plea to appear: ``For we went
astray because Thou didst not appear.''

Augustine hears the psalm in the light of the Passion as well. ``Stir up thy
might, and come to save us'': the one who seemed to have no power when it was
said, ``If Son of God He is, let Him come down from the Cross'', is asked to
show his strength: ``Having been crucified of weakness, rise Thou in power:
`Stir up Thy power, and come Thou, to save us''' (79.2). And on the second
strophe, when the psalm asks God to perfect the vine his right hand planted,
``and upon the Son of man whom thou hast confirmed for thyself'', Augustine
breaks off his explanation: ``What can be more evident? Why do ye still
expect, that we should still explain to you in discourse, and should we not
rather cry out with you in admiration\ldots?'' The Son of man is the
foundation Paul names: ``For other foundation no one is able to lay, except
that which is laid, which is Christ Jesus'' (79.9). Robert Bellarmine, in his
explanation of the psalms, reports the same reading as common ground: ``All agree that this passage refers literally
to Christ'', because the vine was never made perfect until he came
(\work{Explanatio in Psalmos} 79, expl.~15). The son of man stands at the end
of v.~16 in the Latin and in the American psalm, as at the end of v.~18; a
translation made directly from the Hebrew reads ``branch'' at v.~16.

\subsection{Augustine: Christ given, and still to be given}

The Alleluia verse asks, ``Shew us, O Lord, thy mercy; and grant us thy
salvation'' (Ps 85 (84):8). Augustine, reaching this verse in a sermon on the
psalm, reminds his hearers that this is what they sang (\latin{Hoc est
quod cantavimus}) and says what the salvation is:
\latin{salutare tuum, Christum tuum}, ``Thy salvation, that is, Thy Christ.''
Bellarmine paraphrases the petition as ``Send us your Son for a Savior''
(\work{Explanatio in Psalmos} 84, expl.~7). Five verses later the same psalm
has the verse the Communion Antiphon sings, ``the Lord will give goodness:
and our earth shall yield her fruit'' (Ps 85 (84):13), and Bellarmine reads
it of the Incarnation: the Lord sends his Spirit to overshadow a virgin, and
``our land, which was never ploughed nor sown, and was altogether an
untouched virgin, will yield her fruit'' (expl.~12).

If the cry was answered at Bethlehem, why does the Church pray it again?
On the verse he has just explained, Augustine says why the Church still asks
for a Christ already given. \latin{Jam quidem
dedit nobis Christum suum: adhuc illi tamen dicamus, Da nobis Christum
tuum}: ``He has already given us his Christ; yet let us still say to him,
Give us thy Christ.'' We say it, he explains, as we say ``Give us this day our
daily bread'', and our bread is the one who said, ``I am the living bread
which came down from heaven.'' Then he gives a second reason: \latin{Dedit
enim nobis Christum, sed hominem: quem nobis dedit hominem, eum ipsum nobis
daturus est Deum}, ``For he gave us Christ, but as man; the one he gave us as
man, the same he will give us as God.'' He gave him to men as men could
receive him, for no man could receive Christ as God. We are sons of God
already, Augustine goes on, but by faith and in hope, not yet in fact, and
``it doth not yet appear what we shall be'' (1~Jn 3:2). At the end of the
section he returns to the verse: ``grant us Thy Christ, that we may know Thy
Christ, see Thy Christ; not as the Jews saw Him and crucified Him, but as the
Angels see Him, and rejoice'' (84.9; §6 in the Oxford translation, which
passes over the sentences on Christ already given and still to be given; the
English of those sentences is the study's, from the Latin).

\subsection{Cyril: not one coming but two}

Cyril of Jerusalem, instructing the candidates for baptism in his city,
teaches the two comings for which the Preface gives thanks: ``We preach not one
advent only of Christ, but a second also, far more glorious than the former.
For the former gave a view of His patience; but the latter brings with it the
crown of a divine kingdom.'' The descents, he says, are two: ``one, the
unobserved, like rain on a fleece; and a second His open coming, which is to
be''; and so, ``We rest not then upon His first advent only, but look also
for His second'' (\work{Catechetical Lectures} 15.1).

Nothing at Bethlehem tore the sky or melted the mountains. Taken with Cyril's
distinction, Isaiah's image fits the open coming better than the hidden one,
while for Jerome the hidden coming already answered the cry. The two need not
be set against each other: what the prayer asked to see, the sky torn open,
is still to come. The Preface names the two
comings: the first \latin{in humilitate carnis assumptae},
in the lowliness of the flesh taken on, which opened the way of salvation;
the second \latin{in suae gloria maiestatis}, in the glory of his majesty,
when those who now wait will receive openly what was promised. The solemn
blessing, where it is used, names them together in its last invocation:
those who rejoice in the coming in the flesh ask for eternal life at the
coming in majesty.

Paul names both comings in one sentence: he thanks God for gifts already
given in Christ, ``waiting for the manifestation of our Lord Jesus Christ''
(1~Cor 1:7). John Chrysostom presses the word: ```Revelation' is his word;
implying that although He be not seen, yet He is, and is present even now,
and then shall appear'', and so ``there is need of patience''
(\work{Homilies on First Corinthians} 2). Aquinas says of this waiting that
those who wait for his revelation are blessed now in hope (\work{Super I ad
Corinthios} c.~1, lect.~1). The Collect asks for the will to go out to meet
Christ as he comes (\latin{venienti}). The Prayer over the Offerings
asks that what devotion does in this passing time become the reward of
eternal redemption.

In the Gospel the master of the house has already been there and has gone
abroad, and Bede, using Gregory the Great's words on the traveller of
Matthew~25, reads him as Christ ascended. Isaiah's people ask God to come
down; the disciples are told to stay awake for one who will come back.

\subsection{Where this interpretation strains}

Jerome and Aquinas disagree over the verb on which Jerome's answer rests. For
Jerome the past tense is the answer; for Aquinas, in the sense he prefers, it
is a memory, the ground on which the people dare to ask. Aquinas does not
deny Jerome's reading: he keeps the Incarnation as a sense of the words and
adds the judgment, and so reads the one verb of three descents: the descent
that freed Israel from Egypt, the descent into the womb that Jerome saw
fulfilled, and the descent to judge that is still to come. None of the
Fathers cited sets Isaiah's cry beside Mark's parable.

\subsection{The four senses of this interpretation}

\begin{fourSenses}
\item[Literal.] A people that calls the Lord its Father and Redeemer, its
sanctuary trodden down, prays that God tear open the sky and come down as he
came down to free Israel from Egypt (Aquinas's literal sense); the psalmist
asks the Shepherd of Israel to stir up his might and come to save the ruined
vine.
\item[Allegorical.] The descent is the Incarnation: having prayed ``come
down'', the people say ``thou didst come down'', and the Word was made flesh
(Jerome; the second sense Aquinas lists). The son of man on whom God's hand rests is
Christ (Augustine, Bellarmine), and the salvation the Alleluia asks for is
Christ (Augustine). The first descent was unseen, like rain on a fleece
(Cyril).
\item[Moral.] Each believer prays the cry as his own and lifts up his soul
(the Entrance Antiphon); he asks to be turned and shown God's face (the psalm
refrain), and for the will to go out to meet Christ (the Collect).
\item[Anagogical.] The second descent is in glory, open, with the crown of a
divine kingdom (Cyril; the first sense Aquinas lists, the past put for the
future); the Christ once given as man will be given as God, and seen as the
angels see him (Augustine). The Preface says that at his coming in glory
those who now wait will receive openly what they were promised, and the
blessing asks eternal life at the coming in majesty. Until then, the Prayer after
Communion says, God teaches those who receive the mysteries, among passing
things, to love what is heavenly and hold to what lasts.
\end{fourSenses}
